\documentclass[12pt, twocolumn, landscape, a4paper]{article}

\usepackage[margin=1in]{geometry}
\usepackage[utf8]{inputenc}
\usepackage[spanish,es-noshorthands]{babel}
\usepackage{amssymb, amsmath, amsbsy}
%\usepackage{verbatim}
\usepackage{mathpazo}
%\usepackage{tasks}
%\usepackage{enumerate}
\usepackage{graphicx} % \scalebox{1.4}{$$}
\usepackage{tikz}
\usetikzlibrary {arrows.meta} 
% \usetikzlibrary{angles,quotes}
\usepackage[pdftex]{hyperref}
\hypersetup{pdftitle={Ejemplo de vector-fuerza},pdfauthor={Luis Carrillo Gutiérrez}}

\renewcommand{\familydefault}{\sfdefault}
\pagestyle{empty} 

\begin{document}

\begin{itemize}
\item{Ejemplo 1 --- Una persona ejerce una fuerza $F_{\,1} = 4.5\,Kg{-}f$, mientras otra (persona) aplica una fuerza $F_{\,2} = 4\,Kg{-}f$ con un ángulo de $37^{\,\circ}$ ambas respecto al sur.

\begin{tikzpicture}
	\begin{scope}[gray,<->]
	\draw (-2.5,0) -- (2.5,0);
	\draw (0,-5) -- (0,2.5);
	\end{scope}
	\draw (-.35, 2.5) node {N};
	\draw (-.35, -5) node {S};
	\draw (2.5, -.35) node {E};
	\draw (-2.5, -.35) node {O};
	\begin{scope}[-{Latex[length=3.1mm]}, line width=1.3]
	\draw (0, 0) -- (0, -4.5);
	\draw (0, 0) -- (-1.5, -4);
	\end{scope}
	\draw (-2, -3.6) node {$F_{\,2}$};
	\draw (.5, -4) node {$F_{\,1}$};

	% \draw (-.25, 0) arc (265:270:2);
\end{tikzpicture}
}
\end{itemize}

\end{document}